% Pursuing Stacks (cotes 134-2 to 134-8): archive cote 134-4, site page 52 (PDF page 53), typescript page 140.
% typescript. Draft compared: ps3.tex:2635-2702.
% AI-transcribed by Claude Opus 5.5 (claude-opus-5-5, adaptive thinking, effort high), Message Batches API, batch ps2 (msgbatch_01NVwr1Yuo8qJf5rykeGBXPK), 2026-11-12; not yet checked by eye.
\scanpage{134-4/52}
\mspage{140}
fied if I is connected\hand{,} \inserted{In fact}\xout{\illegible{1}} (i') is satisfied iff $\hand{\Sigma}_h$ admits a gene\xout{r}
rating subset $\hand{\Sigma}_h^o$ made up with \uline{connected} intervals.

(ii') Any contractible ob\over{j}{j}ect I is \xout{co}connected \inserted{and}\hand{(}hence 0-connected).

\xout{\illegible{1}} Conversely\inserted{(cf.\ page 130),} if condition e) holds i.e. $\hand{\Pi}_o$ commutes to pro-
ducts, \xout{then} and if moreover h is definable in terms of $\hand{\Sigma}_h$ i.e. comes
from a weak homotopy interval structure, then (i') implies 0-admissibili-
ty . If M admits even a generating family of contractible weak homotopy
intervals, namely if h comes from a contractibility structure on M, then
(ii') equally implies \xout{compatibil} 0-admissibility.

\uline{Remarks}.1) We can generalise these converse statements, by dropping the
condition e) on M, but demanding instead that the connected component
invol\xout{v}ved (namely I itself in case (ii')) is not only 0-connected, but
even ``0-\uline{connected} \uline{over} $e_M$'' , which just means that it's products by any
\inserted{0-}connected object of M is again 0-connected. \inserted[typed]{In the case (ii')} \over{T}{t}his condition is equally \uline{ne-%
\xout{In case (ii')}
cessary} for admissibility. (The corresponding statements could have been
made in the general context of page 130 of course ...)

2) The name of 0-admissibility suggests there may exist corresponding
``higher'' notions of n-admissibility for h, \inserted[typed]{where n is any natural integer.} I \uline{do} see \xout{much} a natural candi-%
date, namely whenever we have a functor \xout{fx} $\hand{\Pi}_n$ from M to n-truncated
homotopy types (as is the case, say, when M is either a topos - we
then rather get \uline{pro}homotopy types - or a modelizer). But it would seem
that in all cases of geometrical significance, and when moreover h is
defined by a weak homotopy interval structure, that 0-admissibility
implies \inserted[typed]{already} n-admissibility for any n.

Due to the existence of \xout{arbitrary} sup for an arbitrary subset\hand{,} in
each of the ordered sets Hom$_{\hand{i}}$(M) ($i\hand{\in\{0,3,4\}}$) of homotopy structures in
M \hand{(} either unqualified, or weak homotopy interval structures, or con-%
tractibility structures\hand{)}, it follows that for each of these three types
of homotopy structures, there is a \xout{"}widest one $h_i$ among all those of
this type, which are 0-admissible. We are interested here, because of
topological motivations, by the case of Hom$_{\hand{3}}$ , namely weak homotopy in-%
terval structures. We call the corresponding homotopy structure $h=h_{\hand{3}}$
the \uline{canonical} \uline{homotopy} \uline{structure} of M. In case M is totally 0-connected
(by which we mean condition e)), the weak homotopy intervals for this
structure are just those intervals for which $\hand{\delta}_o$ , $\hand{\delta}_1$ correspond to the
same connected component of I - and we get a generating set of weak homo-%
topy intervals, by just taking \uline{all connected intervals}. \deleted{If M is not}
\hand{\uline{Remark} If M is not}

\margin{52}
